\documentclass[11pt]{article}
\usepackage{amsmath,amssymb}
\usepackage[margin=1in]{geometry}
\usepackage{booktabs}

\newcommand{\BB}{\mathsf{B}}
\newcommand{\WW}{W}

\title{Disproof of TMC Conjecture 00000000433:\\
the Case $\BB_2$ of the Denominator-Terms Conjecture}
\author{TLMC falsification pipeline submission}
\date{October 2, 2026}

\begin{document}
\maketitle

\begin{abstract}
Conjecture 00000000433 asserts that if $T(\mathrm{type})$ denotes the number
of monomial terms in Macdonald's denominator identity, then
$T(\BB_n)=T(C_n)=n$, $T(D_n)=2n$, $T(G_2)=2$, $T(F_4)=4$, $T(E_6)=6$,
$T(E_7)=7$, $T(E_8)=8$, and that the number of terms always equals the
multiplicity of the exponent $1$. We disprove the conjecture at $\BB_2$:
a direct computation of the Weyl-group orbit of $\rho$ shows the denominator
identity of $\BB_2$ has exactly $8$ monomial terms, i.e.\ $T(\BB_2)=8\neq 2=n$.
The failure is general: $T(\mathrm{type})=|\WW(\mathrm{type})|$ for every
root system in the list, so all displayed values are false.
\end{abstract}

\section{The conjecture}

Macdonald's denominator identity expresses the denominator of an affine Lie
algebra as a linear combination of monomials; let $T(\mathrm{type})$ be its
number of terms. The conjecture claims $T(\BB_n)=T(C_n)=n$, $T(D_n)=2n$,
$T(G_2)=2$, $T(F_4)=4$, $T(E_6)=6$, $T(E_7)=7$, $T(E_8)=8$, and that the
number of terms always equals the multiplicity of exponent $1$.

\section{The attack: $\BB_2$}

The finite-type right-hand side of Macdonald's denominator identity is the
Weyl denominator sum
\[
\prod_{\alpha\in\Delta^+}\bigl(1-e^{-\alpha}\bigr)
 \;=\; \sum_{w\in\WW}\varepsilon(w)\,e^{\,w(\rho)-\rho},
\qquad
\rho=\tfrac12\sum_{\alpha\in\Delta^+}\alpha .
\]
Coefficients are $\pm1$, so the number of monomial terms equals the number of
\emph{distinct} values of $w(\rho)-\rho$ as $w$ ranges over $\WW$.

Realize $\BB_2$ orthonormally: $\Delta=\{\pm e_1,\pm e_2,\pm e_1\pm e_2\}$,
$\Delta^+=\{e_1,\ e_2,\ e_1-e_2,\ e_1+e_2\}$, so
\[
\rho=\tfrac12\bigl(e_1+e_2+(e_1-e_2)+(e_1+e_2)\bigr)=\bigl(\tfrac32,\tfrac12\bigr),
\]
and $\WW(\BB_2)$ is the group of the $8$ signed coordinate permutations
$x\mapsto(\pm x_{\sigma(1)},\pm x_{\sigma(2)})$.

\begin{table}[h]
\centering
\begin{tabular}{lc}
\toprule
$w\in\WW(\BB_2)$ & $w(\rho)-\rho$ \\
\midrule
$x\mapsto(x_1,x_2)$            & $(0,0)$   \\
$x\mapsto(-x_1,x_2)$           & $(-3,0)$  \\
$x\mapsto(x_1,-x_2)$           & $(0,-1)$  \\
$x\mapsto(-x_1,-x_2)$          & $(-3,-1)$ \\
$x\mapsto(x_2,x_1)$            & $(-1,1)$  \\
$x\mapsto(-x_2,x_1)$           & $(-2,1)$  \\
$x\mapsto(x_2,-x_1)$           & $(-1,-2)$ \\
$x\mapsto(-x_2,-x_1)$          & $(-2,-2)$ \\
\bottomrule
\end{tabular}
\caption{The eight exponents $w(\rho)-\rho$ for $\BB_2$.}
\end{table}

All eight values are pairwise distinct. Indeed $\rho$ is strictly dominant,
hence has trivial stabilizer, so $w(\rho)=w'(\rho)$ implies $w=w'$; no two
terms of the sum share a monomial and no cancellation can occur. Therefore
\[
T(\BB_2)=8\;\neq\;2\;=\;n,
\]
which falsifies the conjecture. As a byproduct, the companion clause fails
too: the exponents of $\BB_2$ are $\{1,3\}$, so the multiplicity of exponent
$1$ is $1\neq 8=T(\BB_2)$.

\paragraph{Independent cross-check.}
In simple-root coordinates ($\alpha_1$ short; Cartan matrix
$\bigl(\begin{smallmatrix}2&-2\\-1&2\end{smallmatrix}\bigr)$; reflections
$s_1(a,b)=(-a+2b,\,b)$, $s_2(a,b)=(a,\,a-b)$; positive roots
$\alpha_1,\alpha_2,\alpha_1{+}\alpha_2,2\alpha_1{+}\alpha_2$; take the
integral multiple $2\rho=(4,3)$, harmless since scaling by $2$ is a
bijection), the eight values $w(2\rho)-2\rho$ are
\[
(0,0),\,(-2,0),\,(0,-2),\,(-6,-2),\,(-2,-4),\,(-8,-4),\,(-6,-6),\,(-8,-6),
\]
again pairwise distinct, with sanity checks $w_0(2\rho)=-2\rho$ and
$\sum_w (w(2\rho)-2\rho) = -8\,(4,3)$. This also reproduces the pipeline
verdict list entry-for-entry in the orthonormal realization.

\section{Generality of the failure}

Since $\rho$ is strictly dominant for every root system, its stabilizer is
trivial and the $|\WW|$ exponents $w(\rho)-\rho$ are pairwise distinct; hence
$T(\mathrm{type})=|\WW(\mathrm{type})|$ throughout:

\begin{center}
\begin{tabular}{lcc}
\toprule
type & conjectured $T$ & actual $T=|\WW|$ \\
\midrule
$\BB_1$ ($=A_1=C_1$) & $1$ & $2$ \\
$\BB_n,\ C_n$        & $n$ & $2^n n!$ \\
$D_n$                & $2n$ & $2^{n-1} n!$ \\
$G_2$                & $2$ & $12$ \\
$F_4$                & $4$ & $1152$ \\
$E_6$                & $6$ & $51840$ \\
$E_7$                & $7$ & $2903040$ \\
$E_8$                & $8$ & $696729600$ \\
\bottomrule
\end{tabular}
\end{center}

Even $n=1$ fails ($T(\BB_1)=2\neq1$), so there is no boundary case in which
the conjecture holds.

\section{Machine verification}

The accompanying \texttt{reproduce.py} recomputes both realizations in exact
rational arithmetic and asserts the table above entry-for-entry. The
directory \texttt{lean4/} contains a core-Lean (no Mathlib) formalization:
the eight exponents are produced by the eight explicit signed permutations,
and their pairwise distinctness, the count $8$, and the negation of
$T(\BB_2)=2$ are proved by \texttt{rfl}/\texttt{decide} alone. Every theorem
is axiom-free: \texttt{\#print axioms} reports ``does not depend on any
axioms'' for each, with no \texttt{sorry} and no \texttt{native\_decide}.

\section{Conclusion}

\textbf{Verdict: FALSE.} Conjecture 00000000433 fails at $\BB_2$:
$T(\BB_2)=8\neq 2$, and in fact $T(\mathrm{type})=|\WW(\mathrm{type})|$
for every root system named by the conjecture.

\end{document}
